\section{The Polite Art of Not Being Harmed}

The young believe that the highest form of defence is a wall of flame. The middle-aged believe it is a well- placed parry. The old know that the highest form of defence is not being where the harm lands. The Weave agrees with the old.

These spells do not strike. They do not burn. They do not avenge. They simply \emph{interpose}. A layer of still air between your ribs and the dagger. A whisper of misdirection that turns a sword's edge a finger's width from your throat. A ward that asks, politely, \emph{``Would you mind terribly not?''}

The Weave is not a shield. It is a negotiation. These spells are the terms.

\subsection{On the Nature of Defensive Magic}

A defensive spell does not eliminate danger. It \emph{repositions} it. A [WARD] does not make a door unbreakable; it makes the door inconvenient to break. A [DEFEND] does not make you invulnerable; it makes the attacker's blade hesitate at the moment of commitment.

The tags in this chapter are:

\begin{itemize}
    \item \textbf{[WARD]} — A threshold effect. Something may not cross. The Weave enforces this as a suggestion, not a law, but the suggestion is strong.
    \item \textbf{[DEFEND]} — A personal barrier. The Weave thickens the air, hardens the skin, or simply distracts the eye. Harm becomes Fatigue. Steel becomes discomfort.
    \item \textbf{[PROTECT]} — An interposition. The Weave places the caster between a threat and a target. Dangerous. Often heroic. Usually foolish.
    \item \textbf{[COUNTER]} — A response. The Weave unravels another spell as it forms. Timing is everything. A moment too late, and you are swallowing fire.
    \item \textbf{[REFLECT]} — A reversal. The Weave catches an incoming effect and sends it back toward its source. The Weave does this reluctantly. It dislikes arguments.
\end{itemize}

\subsection{Common Defensive Spells}

\noindent
\begin{longtable}{p{0.2\textwidth} p{0.25\textwidth} p{0.45\textwidth}}
\caption{Defensive Spells of the Wary Caster}\\
\midrule
\textbf{Spell} & \textbf{Tags} & \textbf{What It Does (in practice)}\\
\midrule
\endfirsthead
\multicolumn{3}{c}{\tablename\ \thetable{} -- continued}\\
\midrule
\textbf{Spell} & \textbf{Tags} & \textbf{What It Does (in practice)}\\
\midrule
\endhead
Still Air & [DEFEND] & A cushion of motionless air settles around the caster. The first physical attack against them within the next few heartbeats loses much of its force. A sword that would have opened a vein leaves a bruise. A bruise is survivable. \\
Warded Threshold & [WARD] & The caster touches a door, a window, or a bridge. For a short while, the Weave discourages crossing. A guard will find a reason to patrol elsewhere. A pursuer will hesitate. The effect is psychological, not physical. The Weave is polite. \\
Intercept & [PROTECT][NEAR] & When an ally within Near range is attacked, the caster may interpose themselves. The harm intended for the ally lands on the caster instead. The Weave permits this transfer but does not mitigate it. Use sparingly. \\
Unweave & [COUNTER] & The caster watches an enemy cast a spell and attempts to speak the counter-grammar before the Weave answers. On success, the enemy's spell dissipates harmlessly. On failure, both spells may misfire. The Weave dislikes being interrupted. \\
Shield of Reflection & [WARD][REFLECT] & A shimmering disc of barely- visible force hovers before the caster. The first projectile—arrow, bolt, thrown knife—that strikes the disc is turned back toward its source. The returned projectile is less accurate. The Weave does not aim. \\
Circle of Salt & [WARD][AREA] & The caster draws a line of salt, chalk, or ash in a circle around themselves and allies. For the duration of a scene, creatures of spirit—ghosts, wraiths, some fae—cannot cross the line. Mortals may step over with a shrug. The Weave distinguishes between kinds of hunger. \\
Phantom Parry & [DEFEND][STRIKE] & The Weave conjures a momentary replica of a blade the caster is holding. The replica intercepts one melee attack, then vanishes. The caster does not need to move. The Weave does the work. The caster must still breathe. \\
\hline
\end{longtable}

\subsection{On Failure, Protective}

Defensive spells that fail do not merely fail. They \emph{invert}. The Weave's protection becomes a vulnerability.

\begin{tabular}{p{0.25\textwidth} p{0.7\textwidth}}
\textbf{Still Air (Miss)} & The cushion of air becomes a vacuum. The caster's breath catches. Mark one Fatigue. The attack lands without mitigation. \\
\textbf{Warded Threshold (Partial)} & The ward holds, but the Weave's suggestion is weak. A determined pursuer may overcome it with a moment's effort. The GM gains a Story Beat. \\
\textbf{Intercept (Partial)} & The caster interposes, but the transfer is incomplete. Both the caster and the ally take reduced harm. Neither is safe. Both are angry. \\
\textbf{Unweave (Miss)} & The counter-grammar fails. Worse, it amplifies. The enemy's spell resolves with greater force. The GM gains two Story Beats. \\
\textbf{Phantom Parry (Miss)} & The replica blade appears, but it is oriented incorrectly. The caster's own defence is compromised. The next attack against them has Dominant Position. \\
\end{tabular}

\subsection{A Lesson from the Arena}

A student once asked me, \emph{``Why would I cast a defensive spell when I could simply not be there?''}

She was correct. The best defence is absence. But you cannot always be absent. Sometimes you must stand where the harm is, because someone else cannot run.

I recall a graduate of the Class of Eight Hundred and Forty- Three—a woman named Calypso, now a captain in the Oshiiran lancers. She stood between a charging elephant and a fallen comrade. She had no time to move. She had time to cast [DEFEND][PROTECT].

She rolled a partial success. The Weave interposed. The elephant struck. She took the harm meant for the fallen soldier and walked with a limp for a month. The soldier survived. Calypso called it a fair trade.

The Weave does not promise fairness. It promises \emph{exchange}.

\subsection{When to Defend}

Cast a defensive spell when the alternative is bleeding.

Cast it when you are protecting someone who cannot protect themselves.

Cast it when you have failed to be absent.

Do not cast it when you are safe. The Weave respects economy. A ward that is never tested is a lecture no one attended.

\subsection{A Final Observation on Wards}

The most effective ward I have ever cast was not a spell. It was a lie. I told a Chain- Lantern inquisitor that the room behind me was full of reliquary guards. He believed me. He walked away.

The Weave is a tool. Deception is another. Learn both.

\textit{--- Lysandra}
